% Generated draft: requires mathematical and source review.
\documentclass[11pt]{article}
\usepackage[margin=25mm]{geometry}
\usepackage{fontspec}
\setmainfont{Latin Modern Roman}
\usepackage{amsmath,amssymb,mathtools,unicode-math}
\setmathfont{Latin Modern Math}
\usepackage{enumitem,needspace,xurl,hyperref}
\setlength{\parindent}{0pt}
\setlength{\parskip}{5pt}
\setlength{\emergencystretch}{3em}
\newcommand{\N}{\mathbb{N}}
\newcommand{\Z}{\mathbb{Z}}
\newcommand{\Q}{\mathbb{Q}}
\newcommand{\R}{\mathbb{R}}
\newcommand{\C}{\mathbb{C}}
\newcommand{\F}{\mathbb{F}}
\newcommand{\A}{\mathbb{A}}
\newcommand{\PP}{\mathbb{P}}
\newcommand{\E}{\mathbb{E}}
\newcommand{\eps}{\varepsilon}
\newcommand{\dd}{\,\mathrm{d}}
\begin{document}
\section*{UnsolvedMath ID: 30002151}
\textbf{Source problem code:} catalogue entry
\subsection{Definitions and mathematical statement}
Let \(k\) be a number field, let \(\bar k\) be an algebraic closure, and put \(G_k=\operatorname{Gal}(\bar k/k)\). Let \(X/k\) be a smooth, geometrically connected curve, and choose a geometric base point \(\bar x\) on \(X_{\bar k}\). Throughout, fundamental groups are profinite étale fundamental groups. Assume that \(\pi_1(X_{\bar k},\bar x)\) is nonabelian. There is an exact sequence
\[
1\longrightarrow \pi_1(X_{\bar k},\bar x)\longrightarrow \pi_1(X,\bar x)\xrightarrow{p}G_k\longrightarrow 1.
\]
A section is a continuous group homomorphism \(s:G_k\to\pi_1(X,\bar x)\) satisfying \(p\circ s=\operatorname{id}_{G_k}\). Two sections are equivalent if they are conjugate by an element of \(\pi_1(X_{\bar k},\bar x)\). Write \(\operatorname{Sec}(X/k)\) for the set of equivalence classes.

To make the word \emph{noncuspidal} precise, use the following standard boundary interpretation, which is not explicitly defined in the supplied record. Let \(\overline X\) be the smooth projective compactification of \(X\), and let \(B=\overline X\setminus X\) be its reduced boundary. For \(c\in B(k)\), choose a decomposition subgroup \(D_c\subseteq\pi_1(X,\bar x)\) associated with a prolongation of the valuation at \(c\) to the maximal extension of \(k(X)\) unramified over \(X\). Call a section class \([s]\) cuspidal if, for some \(c\in B(k)\) and some representative \(s'\) of \([s]\),
\[
s'(G_k)\subseteq D_c.
\]
The existence condition is independent of the choices of decomposition subgroups. Let \(\operatorname{Sec}_{\mathrm{nc}}(X/k)\) denote the classes that are not cuspidal. If \(X\) is projective, its boundary is empty and every section class is noncuspidal.

Each \(a\in X(k)\), viewed as a morphism \(\operatorname{Spec}k\to X\), induces a section up to geometric conjugacy. Thus there is a point-to-section map
\[
\kappa_X:X(k)\longrightarrow\operatorname{Sec}(X/k),\qquad a\longmapsto[s_a].
\]
\textbf{Statement requested by the supplied record (Grothendieck's section conjecture).} For every number field \(k\) and every smooth, geometrically connected curve \(X/k\) with nonabelian geometric étale fundamental group, is \(\kappa_X\) injective with image exactly \(\operatorname{Sec}_{\mathrm{nc}}(X/k)\)? Equivalently, does it induce a bijection
\[
X(k)\xrightarrow{\sim}\operatorname{Sec}_{\mathrm{nc}}(X/k)?
\]
The supplied record labels this problem open; that label is supplied curation, not an independently verified status assessment. The boundary interpretation above requires comparison with the cited report during mathematical review.
\subsection{Short English statement}
For every number field and every smooth, geometrically connected curve with nonabelian geometric étale fundamental group, the conjecture asks whether rational points correspond bijectively to continuous sections of the arithmetic fundamental-group sequence, up to geometric conjugacy, after excluding sections arising at boundary points. The supplied source does not explicitly define the boundary exclusion; the draft states the standard interpretation for review.
\subsection{Sources}
\begin{itemize}
\item Local UnsolvedMath registry, supplied complete source record: permanent ID \(30002151\), exact source problem code \texttt{catalogue entry}, title \emph{Grothendieck Section Conjecture for Hyperbolic Curves}, collection \emph{Oberwolfach Reports Open Problems}. The conjecture is attributed to Grothendieck; no individual proposer of this report entry is supplied. The record gives proposed year \(2012\).
\item Supplied report citation: \emph{Arithmetic Geometry (2013)}, Oberwolfach Reports, DOI \texttt{10.4171/owr/2012/38}, pp.~2335--2388. \url{https://doi.org/10.4171/owr/2012/38}. The background describes Oberwolfach Reports as documenting mathematical workshops and their open-problem discussions. The citation's title year \(2013\) differs from the \(2012\) appearing in the DOI and record year; this discrepancy is retained without resolution.
\item Supplied literature source: \emph{Stix, section conjecture for curves of index one}. \url{https://jtnb.centre-mersenne.org/articles/10.5802/jtnb.1362/}. No additional bibliographic details are supplied.
\item Literature assessment checked at \texttt{2026-08-21}: the supplied curation labels the problem open and describes the general section conjecture for hyperbolic curves as remaining wide open, with special cases and local or conditional results rather than a general proof. It lists the report and the Stix reference above as sources checked. This is supplied dated curation derived from the OWR material, not independently verified here; the status and cited claims require independent verification.
\item The supplied statement-status label is \texttt{corrected\_verified}. Its verification note says that the official report states this formulation and all its curve hypotheses. This is a supplied verification claim; the report has not been independently inspected for this draft.
\end{itemize}
\end{document}
